\documentclass[10pt,a4paper]{amsart}
\usepackage[margin=19mm]{geometry}
\usepackage{amsmath,amssymb,mathtools}
\usepackage[scr=rsfs,cal=euler]{mathalfa}
\usepackage{xfrac}
\usepackage[unicode,hidelinks,bookmarks=false]{hyperref}
\newcommand{\ie}{i.e.\ }
\newcommand{\cf}{cf.\ }
\newcommand{\resp}{respectively\ }
\newcommand{\Subsection}{\subsection}
\providecommand{\nobreakdash}{\nobreak\hbox{-}}
\setlength{\emergencystretch}{2em}
\multlinegap0pt
\usepackage{bm}
\usepackage{bbm}
\usepackage{dsfont}
\usepackage{xspace}
\usepackage{cancel}
\usepackage{mathtools}
\usepackage{amsthm}
\newtheorem{theorem}{Theorem}
\title[Applications of sequential spaces to Smyth power spaces and ]{Applications of sequential spaces to Smyth power spaces and sober spaces}
\author{Zhengmao He}
\date{Source: arXiv:2609.27775v1 (CC BY 4.0)}
\begin{document}
\maketitle

\noindent\emph{Source and licence.} Applications of sequential spaces to Smyth power spaces and sober spaces. Version arXiv:2609.27775v1 (CC BY 4.0), distributed under the Creative Commons Attribution 4.0 International licence, as recorded in the arXiv licence metadata for this exact version and shown on the abstract page https://arxiv.org/abs/2609.27775v1. Full rights notice: licenses/arxiv-2609.27775-CC-BY-4.0.txt.\par\smallskip

\noindent\textit{Editorial scope.} Counted unit: the Theorem whose source label is \texttt{None} in the article (source0.tex lines 608--609, source label \texttt{None}), delivered together with the article's own complete original proof of it (source0.tex lines 612--640, one \texttt{proof} environment of 29 lines). No mathematics is written, repaired or re-derived: statement and proof are the article's own text line for line. Required context retained uncounted: none. Every definition, display and label of those lines is retained unchanged. Omitted from this package and not redistributed: 1--607, 610--611, 641--850. Nothing inside a retained excerpt is omitted or altered: every retained body is delivered exactly as the article's lines are written, so no omission receipt of any kind accompanies this package. Citations: the retained excerpts carry no citation command at all, so no bibliography entry of the article is transcribed and no reference is redistributed. Reference adaptation: none was needed and none was made; every reference inside the delivered unit resolves inside it, so no unresolved reference or citation is produced. Printed slips kept literal: no printed word, symbol, hypothesis or number of the counted statement or of its proof was altered. Layout and class: the article's own class is \texttt{elsarticle} with packages including enumerate, amssymb, lipsum, geometry, booktabs, amsmath, url, enumitem, graphics, array, dsfont, xy, amsthm, tikz; the wrapper is the lane's pinned \texttt{amsart} editorial wrapper at 10pt on A4, which restates the theorem-like environments the retained text needs and adds the article's own macro definitions that the retained text needs (none), with extra wrapper packages (bm, bbm, dsfont, xspace, cancel, mathtools). The delivered numbering is the unit's own. Assets: no retained span contains a figure environment, a graphics-inclusion command, a PDF-page inclusion, a table or a TikZ picture; no image file is copied into this package, the package declares no asset dependency and no image file is redistributed with it. Licence: version arXiv:2609.27775v1 of this article is distributed under the Creative Commons Attribution 4.0 International licence, as recorded in the arXiv licence metadata for this exact version and shown on the abstract page https://arxiv.org/abs/2609.27775v1. A full rights notice is delivered at licenses/arxiv-2609.27775-CC-BY-4.0.txt. Characters: every retained excerpt is pure ASCII, so the delivered engine, XeLaTeX with its default Latin Modern text font, reports no missing character.
\begin{theorem} {\rm Let $P$ and $Q$ be two posets. If $\Sigma P\times \Sigma Q$ is a sequential  space, then $\Sigma (P\times Q)=\Sigma P \times \Sigma Q$.}
\end{theorem}

\begin{proof}
Set $$\tau=\sigma(P)\times\sigma(Q)
    \quad\text{and}\quad
    \rho=\sigma(P\times Q).$$
It is always the case that $$\sigma(P)\times\sigma(Q)\subseteq\sigma(P\times Q).$$
Conversely, let $A\subseteq P\times Q$ be
$\rho$-closed. Since $(P\times Q,\tau)=\Sigma P\times\Sigma Q$
is sequential, it suffices to show that $A$ is sequentially closed
with respect to $\tau$. Let $\{(p_n,q_n)\}_{n\in\mathbb{N}}$ be a sequence in $A$ and suppose that
the sequence $\{(p_n,q_n)\}_{n\in\mathbb{N}}$ converges to $(p,q)$ in $\Sigma P\times\Sigma Q$. Then $\{p_n\}_{n\in\mathbb{N}}$ ($\{q_{n}\}_{n\in\mathbb{N}}$) converges to $p$($q$) in $\Sigma P$ ($\Sigma Q$). Assume that $(p,q)\notin A$. Put $$T=(P\times Q)\setminus A.$$ Then $T\in\sigma(P\times Q)$ and $(p,q)\in T$.

Define a mapping $f:P\longrightarrow\mathcal O(\Sigma Q)$ by
 $$ \forall\ a\in P, f(a)=\{b\in Q:(a,b)\in T\}.$$
For every $a\in P$, we can check that $f(a)$ is Scott open in $Q$. Furthermore,
$f$ is Scott continuous. Since $(p,q)\in T$, we have $q\in f(p)$. As
$\{q_{n}\}_{n\in\mathbb{N}}$) converges to $q$, there exists $n_{0}$ such
that for every $n_0\leq n$, $q_n\in f(p)$. Let
$$E=\{q\}\cup\{q_n\mid n_0\leq n\}.$$
Then $E$ is compact in $\Sigma Q$, since it consists of a
convergent sequence together with its limit, and $E\subseteq f(p)$. Define
$$\Box E=\{G\in\mathcal O(\Sigma Q)\mid E\subseteq G\}.$$
Then $\Box E$ is Scott open in $\mathcal O(\Sigma Q)$. Since $E\subseteq f(p)$, we have $f(p)\in\Box E$. By the Scott
continuity of $f$, $W=f^{-1}(\Box E)$
is a Scott open neighbourhood of $p$ in $P$. Since $\{p_n\}_{n\in\mathbb{N}}$ converges to $p$,
there exists $n_0\leq n_1$ such that $p_{n}\in W$ for all $n_1\leq n$. In particular, $p_{n_1}\in W$, so $E\subseteq f(p_{n_1})$. Since $q_{n_1}\in E$, it follows that $q_{n_1}\in f(p_{n_1})$ and therefore $p_{n_1},q_{n_1})\in T$. This contradicts $(p_{n_1},q_{n_1})\in A$. Hence $(p,q)\in A$, and therefore $A$ is sequentially closed in
$(P\times Q,\tau)$. Since $\Sigma P\times \Sigma Q$  is sequential, $A$ is
$\tau$-closed. Thus every $\rho$-closed set is $\tau$-closed,
which gives $\rho\subseteq\tau$. Therefore
$\sigma(P\times Q)=\sigma(P)\times\sigma(Q)$. \end{proof}

\end{document}
